\documentclass[10pt,a4paper]{amsart}
\usepackage{amsmath,amssymb,amsthm,graphicx}
% Retained preamble for the Stacks Project chapter divisors.tex (Divisors).
% The chapter file pinned-original-source.tex carries only \input{preamble} on its first line, so no
% preamble text lives inside the chapter itself.  The Stacks Project uses one single shared preamble for
% every chapter; the snapshot stages it verbatim in sources/stacks-chow/preamble.tex.  The block below is
% copied mechanically, line for line, from that file's lines 14-18 (the Xy-pic setup) and lines 154-217
% (the native theorem declarations together with the mathematical macros).  Nothing was added, removed or
% redefined.  The wrapper class and the amsmath/amssymb/amsthm/graphicx base packages are supplied by the
% assembler, so the class/\externaldocument part of the shared preamble is not repeated here.
% For commutative diagrams we use Xy-pic
\usepackage[all]{xy}

% We use 2cell for 2-commutative diagrams.
\xyoption{2cell}
% Theorem environments.
%
\theoremstyle{plain}
\newtheorem{theorem}[subsection]{Theorem}
\newtheorem{proposition}[subsection]{Proposition}
\newtheorem{lemma}[subsection]{Lemma}

\theoremstyle{definition}
\newtheorem{definition}[subsection]{Definition}
\newtheorem{example}[subsection]{Example}
\newtheorem{exercise}[subsection]{Exercise}
\newtheorem{situation}[subsection]{Situation}

\theoremstyle{remark}
\newtheorem{remark}[subsection]{Remark}
\newtheorem{remarks}[subsection]{Remarks}

\numberwithin{equation}{subsection}

% Macros
%
\def\lim{\mathop{\mathrm{lim}}\nolimits}
\def\colim{\mathop{\mathrm{colim}}\nolimits}
\def\Spec{\mathop{\mathrm{Spec}}}
\def\Hom{\mathop{\mathrm{Hom}}\nolimits}
\def\Ext{\mathop{\mathrm{Ext}}\nolimits}
\def\SheafHom{\mathop{\mathcal{H}\!\mathit{om}}\nolimits}
\def\SheafExt{\mathop{\mathcal{E}\!\mathit{xt}}\nolimits}
\def\Sch{\mathit{Sch}}
\def\Mor{\mathop{\mathrm{Mor}}\nolimits}
\def\Ob{\mathop{\mathrm{Ob}}\nolimits}
\def\Sh{\mathop{\mathit{Sh}}\nolimits}
\def\NL{\mathop{N\!L}\nolimits}
\def\CH{\mathop{\mathrm{CH}}\nolimits}
\def\proetale{{pro\text{-}\acute{e}tale}}
\def\etale{{\acute{e}tale}}
\def\QCoh{\mathit{QCoh}}
\def\Ker{\mathop{\mathrm{Ker}}}
\def\Im{\mathop{\mathrm{Im}}}
\def\Coker{\mathop{\mathrm{Coker}}}
\def\Coim{\mathop{\mathrm{Coim}}}

% Boxtimes
%
\DeclareMathSymbol{\boxtimes}{\mathbin}{AMSa}{"02}

%
% Macros for moduli stacks/spaces
%
\def\QCohstack{\mathcal{QC}\!\mathit{oh}}
\def\Cohstack{\mathcal{C}\!\mathit{oh}}
\def\Spacesstack{\mathcal{S}\!\mathit{paces}}
\def\Quotfunctor{\mathrm{Quot}}
\def\Hilbfunctor{\mathrm{Hilb}}
\def\Curvesstack{\mathcal{C}\!\mathit{urves}}
\def\Polarizedstack{\mathcal{P}\!\mathit{olarized}}
\def\Complexesstack{\mathcal{C}\!\mathit{omplexes}}
% \Pic is the operator that assigns to X its picard group, usage \Pic(X)
% \Picardstack_{X/B} denotes the Picard stack of X over B
% \Picardfunctor_{X/B} denotes the Picard functor of X over B
\def\Pic{\mathop{\mathrm{Pic}}\nolimits}
\def\Picardstack{\mathcal{P}\!\mathit{ic}}
\def\Picardfunctor{\mathrm{Pic}}
\def\Deformationcategory{\mathcal{D}\!\mathit{ef}}

% --- appended 2026-10-12 (Chief): the project own definitions of the reference/slogan/history
% display environments, copied mechanically from the shared Stacks preamble lines 8-12
% (\usepackage{verbatim} plus the three \newenvironment lines, which the chapter preamble keeps
% verbatim from the project preamble). Nothing added, removed or redefined: these environments
% carry no mathematical content and the project itself typesets them as comments. Without these
% three lines a counted statement whose own text uses them fails with Environment slogan undefined.
\usepackage{verbatim}
\newenvironment{reference}{\comment}{\endcomment}
\newenvironment{slogan}{\comment}{\endcomment}
\newenvironment{history}{\comment}{\endcomment}
\usepackage{cleveref}
\crefname{originallemma}{original lemma}{original lemmas}
\Crefname{originallemma}{Original lemma}{Original lemmas}
\makeatletter
\newlabel{lemma-characterize-OD}{{01X0}{source}{}{sourcepointer-1}{}}
\newlabel{lemma-characterize-OD@cref}{{[originallemma][0][]01X0}{[1][0][]source}{}{}{}}
\makeatother
\crefname{originallemma}{original lemma}{original lemmas}
\Crefname{originallemma}{Original lemma}{Original lemmas}
\makeatletter
\newlabel{lemma-finite-ass}{{05AF}{source}{}{sourcepointer-2}{}}
\newlabel{lemma-finite-ass@cref}{{[originallemma][0][]05AF}{[1][0][]source}{}{}{}}
\makeatother
\crefname{originallemma}{original lemma}{original lemmas}
\Crefname{originallemma}{Original lemma}{Original lemmas}
\makeatletter
\newlabel{lemma-regular-section-associated-points}{{0AYL}{source}{}{sourcepointer-3}{}}
\newlabel{lemma-regular-section-associated-points@cref}{{[originallemma][0][]0AYL}{[1][0][]source}{}{}{}}
\makeatother
\crefname{originallemma}{original lemma}{original lemmas}
\Crefname{originallemma}{Original lemma}{Original lemmas}
\makeatletter
\newlabel{properties-lemma-ample-finite-set-in-principal-affine}{{properties-lemma-ample-finite-set-in-principal-affine}{source}{}{sourcepointer-4}{}}
\newlabel{properties-lemma-ample-finite-set-in-principal-affine@cref}{{[originallemma][0][]properties-lemma-ample-finite-set-in-principal-affine}{[1][0][]source}{}{}{}}
\makeatother
\crefname{originallemma}{original lemma}{original lemmas}
\Crefname{originallemma}{Original lemma}{Original lemmas}
\makeatletter
\newlabel{properties-lemma-invert-s-sections}{{properties-lemma-invert-s-sections}{source}{}{sourcepointer-5}{}}
\newlabel{properties-lemma-invert-s-sections@cref}{{[originallemma][0][]properties-lemma-invert-s-sections}{[1][0][]source}{}{}{}}
\makeatother
\crefname{originallemma}{original lemma}{original lemmas}
\Crefname{originallemma}{Original lemma}{Original lemmas}
\makeatletter
\newlabel{properties-lemma-quasi-affine-invertible-nonvanishing-section}{{properties-lemma-quasi-affine-invertible-nonvanishing-section}{source}{}{sourcepointer-6}{}}
\newlabel{properties-lemma-quasi-affine-invertible-nonvanishing-section@cref}{{[originallemma][0][]properties-lemma-quasi-affine-invertible-nonvanishing-section}{[1][0][]source}{}{}{}}
\makeatother
\title{On a Noetherian scheme with an ample invertible sheaf every invertible module is the difference of two effective Cartier divisors avoiding a prescribed finite set}
\author{The Stacks Project Authors}
\date{stacks-project:divisors.tex:local-snapshot-20261012; GNU FDL 1.2 or later}
\begin{document}
\maketitle
\noindent\textit{Editorial note.} All delivered mathematical text is the author's own, copied from the pinned Stacks Project chapter file pinned-original-source.tex; no statement or proof step was rewritten, reordered or removed.\par\medskip
\noindent\textit{Editorial note.} Following the multi-claim convention for this source, source\_url is the source-level URL https://stacks.math.columbia.edu/ and the source label lemma-quasi-projective-Noetherian-pic-effective-Cartier is carried in claim\_key and here rather than in a per-claim tag URL.\par\medskip
\noindent\textit{Editorial note.} The preamble is the chapter's shared Stacks Project preamble, retained as sources/stacks-divisors/retained-preamble.tex: a line-for-line excerpt of the shared Stacks preamble staged with this snapshot. No macro was invented or redefined, and this card adds no preamble text of its own.\par\medskip
\noindent\textit{Editorial note.} Statement span 3215-3230 and proof span 3232-3274 come from the halfA pool card operations/pool-cards/stacks-divisors-halfA.json and were re-read in the pinned chapter file before this card was written.\par\medskip
\noindent\textit{Editorial note.} The statement carries an author slogan ('On a projective scheme, every line bundle has a regular meromorphic section.'), which is part of the counted statement span and is delivered verbatim.\par\medskip
\noindent\textit{Editorial note.} The statement and proof were checked to contain no citation requiring a bibliography: the only citations are to labelled statements of the project, handled as external navigation pointers.\par\medskip
\noindent\textit{Editorial note.} 6 cross-references occur inside the counted spans; they are listed in external\_references as navigation pointers only.\par\medskip
\noindent\textit{Allowed standard prerequisites (editorial).}\begin{itemize}
\item Associated points of a Noetherian scheme are finite, and a section of an invertible sheaf not vanishing at any associated point is a regular section whose zero scheme is an effective Cartier divisor, with the corresponding invertible sheaf isomorphic to the structure sheaf twisted by that divisor. Used throughout the proof, at lines 3233-3234, 3240-3242 and 3265-3268. Turns a nowhere vanishing condition at finitely many points into regularity and into the description of the invertible sheaf as a divisor class.
\item Existence of a section of a power of an ample invertible sheaf making a prescribed finite set land in an affine open, cited by the author from the properties chapter. Cited at proof lines 3245-3249. Produces the affine open X\_s containing the prescribed points, to which the quasi-affine case is applied.
\item Extendability of sections after multiplication by a power of the defining section, cited by the author from the properties chapter. Cited at proof lines 3253-3257. Produces the global sections of the twisted invertible sheaves whose zero divisors are the announced D and D'.
\end{itemize}\medskip
\ifdefined\hypertarget\hypertarget{sourcepointer-1}{}\fi
\noindent\textit{Source reference.} Original lemma 01X0: focused-seeds/AGP/sources/stacks-snapshot-20261011/tags.txt, source lines 2501--2501. Same-chapter reference used in the retained author text at source line(s) 2902, 3246, 3271. It is supplied only as a navigation pointer so that the retained text compiles; the counted item delivers the author's own statement and the author's complete proof, and no part of the referenced item is copied in or substituted for the counted argument.\par\medskip
\ifdefined\hypertarget\hypertarget{sourcepointer-2}{}\fi
\noindent\textit{Source reference.} Original lemma 05AF: focused-seeds/AGP/sources/stacks-snapshot-20261011/tags.txt, source lines 6848--6848. Same-chapter reference used in the retained author text at source line(s) 3234. It is supplied only as a navigation pointer so that the retained text compiles; the counted item delivers the author's own statement and the author's complete proof, and no part of the referenced item is copied in or substituted for the counted argument.\par\medskip
\ifdefined\hypertarget\hypertarget{sourcepointer-3}{}\fi
\noindent\textit{Source reference.} Original lemma 0AYL: focused-seeds/AGP/sources/stacks-snapshot-20261011/tags.txt, source lines 13794--13794. Same-chapter reference used in the retained author text at source line(s) 3243, 3268. It is supplied only as a navigation pointer so that the retained text compiles; the counted item delivers the author's own statement and the author's complete proof, and no part of the referenced item is copied in or substituted for the counted argument.\par\medskip
\ifdefined\hypertarget\hypertarget{sourcepointer-4}{}\fi
\noindent\textit{Source reference.} Original lemma properties-lemma-ample-finite-set-in-principal-affine: sources/stacks-divisors/pinned-original-source.tex, source lines 3254--3254. Same-chapter reference used in the retained author text at source line(s) 3254, 3765. It is supplied only as a navigation pointer so that the retained text compiles; the counted item delivers the author's own statement and the author's complete proof, and no part of the referenced item is copied in or substituted for the counted argument.\par\medskip
\ifdefined\hypertarget\hypertarget{sourcepointer-5}{}\fi
\noindent\textit{Source reference.} Original lemma properties-lemma-invert-s-sections: sources/stacks-divisors/pinned-original-source.tex, source lines 3261--3261. Same-chapter reference used in the retained author text at source line(s) 3261, 7516. It is supplied only as a navigation pointer so that the retained text compiles; the counted item delivers the author's own statement and the author's complete proof, and no part of the referenced item is copied in or substituted for the counted argument.\par\medskip
\ifdefined\hypertarget\hypertarget{sourcepointer-6}{}\fi
\noindent\textit{Source reference.} Original lemma properties-lemma-quasi-affine-invertible-nonvanishing-section: sources/stacks-divisors/pinned-original-source.tex, source lines 3241--3241. Same-chapter reference used in the retained author text at source line(s) 3241, 3630. It is supplied only as a navigation pointer so that the retained text compiles; the counted item delivers the author's own statement and the author's complete proof, and no part of the referenced item is copied in or substituted for the counted argument.\par\medskip
\begin{lemma}
\label{lemma-quasi-projective-Noetherian-pic-effective-Cartier}
\begin{slogan}
On a projective scheme, every line bundle has a regular meromorphic section.
\end{slogan}
Let $X$ be a Noetherian scheme which has an ample invertible sheaf.
Then every invertible $\mathcal{O}_X$-module is isomorphic to
$$
\mathcal{O}_X(D - D') =
\mathcal{O}_X(D) \otimes_{\mathcal{O}_X} \mathcal{O}_X(D')^{\otimes -1}
$$
for some effective Cartier divisors $D, D'$ in $X$. Moreover, given a
finite subset $E \subset X$ we may choose $D, D'$ such that
$E \cap D = \emptyset$ and $E \cap D' = \emptyset$. If
$X$ is quasi-affine, then we may choose $D' = \emptyset$.
\end{lemma}

\begin{proof}
Let $x_1, \ldots, x_n$ be the associated points of $X$
(Lemma \ref{lemma-finite-ass}).

\medskip\noindent
If $X$ is quasi-affine and $\mathcal{N}$ is any invertible
$\mathcal{O}_X$-module, then we can pick a section $t$ of
$\mathcal{N}$ which does not vanish at any of the points
of $E \cup \{x_1, \ldots, x_n\}$, see Properties, Lemma
\ref{properties-lemma-quasi-affine-invertible-nonvanishing-section}.
Then $t$ is a regular section of $\mathcal{N}$ by
Lemma \ref{lemma-regular-section-associated-points}.
Hence $\mathcal{N} \cong \mathcal{O}_X(D)$ where
$D = Z(t)$ is the effective Cartier divisor corresponding to $t$, see
Lemma \ref{lemma-characterize-OD}. Since $E \cap D = \emptyset$
by construction we are done in this case.

\medskip\noindent
Returning to the general case, let $\mathcal{L}$ be an ample invertible sheaf
on $X$. There exists an $m > 0$ and a section
$s \in \Gamma(X, \mathcal{L}^{\otimes m})$ such that $X_s$
is affine and such that $E \cup \{x_1, \ldots, x_n\} \subset X_s$
(Properties, Lemma \ref{properties-lemma-ample-finite-set-in-principal-affine}).

\medskip\noindent
Let $\mathcal{N}$ be an arbitrary invertible $\mathcal{O}_X$-module.
By the quasi-affine case, we can find a section
$t \in \mathcal{N}(X_s)$ which does not vanish at any point
of $E \cup \{x_1, \ldots, x_n\}$.
By Properties, Lemma \ref{properties-lemma-invert-s-sections}
we see that for some $e \geq 0$ the section $s^e|_{X_s} t$ extends to
a global section $\tau$ of $\mathcal{L}^{\otimes me} \otimes \mathcal{N}$.
Thus both $\mathcal{L}^{\otimes me} \otimes \mathcal{N}$ and
$\mathcal{L}^{\otimes me}$ are invertible sheaves which have global sections
which do not vanish at any point of $E \cup \{x_1, \ldots, x_n\}$.
Thus these are regular sections by
Lemma \ref{lemma-regular-section-associated-points}.
Hence $\mathcal{L}^{\otimes me} \otimes \mathcal{N} \cong \mathcal{O}_X(D)$
and $\mathcal{L}^{\otimes me} \cong \mathcal{O}_X(D')$ for some
effective Cartier divisors $D$ and $D'$, see Lemma \ref{lemma-characterize-OD}.
By construction $E \cap D = \emptyset$ and $E \cap D' = \emptyset$
and the proof is complete.
\end{proof}

\end{document}
